
\documentclass[12pt, a4paper, oneside, romanian]{teza-upb}
\setcounter{secnumdepth}{3}
\setcounter{tocdepth}{3}
\usepackage{babel}
\usepackage{graphicx}
\usepackage{babel}
\usepackage{color}
\usepackage{graphicx}
\graphicspath{ {./figuri/} }
\usepackage[acronym]{glossaries}
\usepackage[
  bookmarks,
  bookmarksopen=true,
  pdftitle={Dizertatie},
  linktocpage]{hyperref}


\singlespacing

%%%%%%%%%%%%%%%%%%%%%%%%%%%%%%%%%%%%%%%%%%%%%%%%%%%%%%%%cod colorat
\usepackage{listings}
\usepackage{color}
 \usepackage{url}  
 \usepackage{float}

\definecolor{codegreen}{rgb}{0,0.6,0}
\definecolor{codegray}{rgb}{0.5,0.5,0.5}
\definecolor{codepurple}{rgb}{0.58,0,0.82}
\definecolor{backcolour}{rgb}{0.95,0.95,0.92}
\usepackage[export]{adjustbox}
\lstdefinestyle{mystyle}{
	language=Matlab,
 %   backgroundcolor=\color{backcolour},   
    commentstyle=\color{codegreen},
    keywordstyle=\color{blue},
    numberstyle=\tiny\color{codegray},
    stringstyle=\color{codepurple},
	basicstyle=\ttfamily\scriptsize,
    breakatwhitespace=false,         
    breaklines=true,                 
    captionpos=b,                    
    keepspaces=true,                 
    numbers=left,                    
    numbersep=5pt,                  
    showspaces=false,                
    showstringspaces=false,
    showtabs=false,                  
    tabsize=2
}
 
\lstset{style=mystyle}
%%%%%%%%%%%%%%%%%%%%%%%%%%%%%%%%%%%%%%%%%%%%%%%


\begin{document}

\author{Ionescu Maria Cătălina}

\title{Studiul și implementarea soluțiilor cloud scalabile pentru aplicații web}


\facultatea{Facultatea de Electronică, Telecomunicații și Tehnologia Informației}
\tiplucrare{licență}
%\tiplucrare{dizertație}
\domeniu{Calculatoare și tehnologia informației}
\catedra{Telecomunicații}
\campus{Leu} 
\program{Ingineria informației}
\titlulobtinut{Inginer}
%\titlulobtinut{Master}
\director{Ş.L.Dr.Ing. George Valentin STOICA} 

\submissionmonth{Iunie} 
\submissionyear{2025} 

\beforepreface
\listoffigures
\listoftables
\abbreviations{ 
HTML = HyperText Markup Language\\
HTTP = Hypertext Transfer Protocol \\
   IoT = Internet of Things\\
   VM1 = Mașină virtuală 1\\
   VM2 = Mașină virtuală 2\\
   

   
  
}


\afterpreface 
\chapter{Introducere}
\section{Prezentarea proiectului}

Cerințele privind performanța și scalabilitatea aplicațiilor web sunt în continuă creștere datorită gamei extinse de dispozitive inteligente și IoT (Internet of Things) folosite de către utilizatori pentru a-și îmbunătăți viața de zi cu zi. Numărul mare de dispozitive care accesează în mod concomitent un site web duce la un volum copleșitor de cereri către servere pentru a accesa informația cerută, ceea ce poate cauza o experiență întârziată sau chiar imposibilă a utilizatorul. Datorită acestor circumstanțe, traficului de date crește, iar aplicațiile web trebuie să facă față la un flux considerabil de informații. La momentul actual, o soluție pentru optimizarea aplicațiilor web expuse unui trafic mare de date, dar și pentru a spori satisfacția utilizatorilor este folosirea soluțiilor scalabile a resurselor software și hardware. 

Scalabilitatea \cite{scal} reprezintă capacitatea unui sistem de a crește sau reduce performanța și costurile unui sistem în funcție de cerințele de procesare. Aceasta este de două tipuri: 
\begin{enumerate}
\item Scalabilitatea resurselor software \cite{scal.soft} se referă la abilitatea unei aplicații de a-și adapta sau menține performanța atunci când este o creștere sau o scădere a volumului de utilizatori, date sau operațiuni. 

\item Scalabilitatea resurselor hardware \cite{scal.hard} reprezintă capacititatea unui sistem de a-și mări capacitatea sau de a-și crește performanța utilizând noi resurse hardware. 
\end{enumerate}
De asemenea, folosirea în paralel a unei bazei de date distribuită reprezintă o soluție modernă pentru asigurarea redundanței datelor și a performanței optime într-o aplicație complexă. O bază de date distribuită \cite{bd} este o bază de date controlată de un sistem de gestiune a bazelor de date (Data Base Management System), în care datele sunt distribuite la mai multe calculatoare interconectate.

\section{Scopul proiectului}
Proiectul \textbf{„Studiul și implementarea soluțiilor cloud scalabile pentru aplicații web”} își propune studierea, implementarea și testarea soluțiilor scalabile software pentru o platformă de rezervări online, utilizând framework-ul Flask Python (pentru dezvoltarea aplicației web), MySQL (sistem de gestionare a bazelor de date relaționale) și Nginx (loadbalancer și server web de tip reverse proxy). Scopul acestui proiect este a urmări și analiza capabilitatea aplicației de a răspunde la un număr mare de cereri simultane, îmbunătățirea performanței, extinderea și reducerea redundanței prin implementarea unei soluții scalabile de tip software care include un load balancer și o bază de date distribuită master-sclave. 

Aplicația simulează o platformă de tip rezervări fiind dezvoltată într-un mediul de lucru alcătuit dintr-o stație de lucru, alături de un mediu de testare, constituit de cele două mașini virtuale instalate pe stația de lucru. Împreună, cele trei mașini fac parte din aceeași rețea IP/LAN. De asemenea, la nivel arhitectural s-a implementat o bază de date distribuită de tip master-slave pentru a optimiza operațiile de citire-scriere.  

Obiectivele specifice acestei lucrării sunt:
\begin{itemize}
\item Dezvoltarea aplicației web folosind framework-ul Flask Python care include rute pentru funcționalități de bază precum ar fi: autentificarea (login, register), managementul utilizatorilor (profile), administrarea rezervărilor (reserve) și pagini administrative (admin).

\item Implementarea arhitecturii scalabile prin intermediul sistemului de load balancing folosind Nginx. Nginx are rolul de a distribui cererile trimise între mai multe servere, astfel crescând performanța și puterea de procesare a aplicației web.

\item Optimizarea performanței bazei de date cu un server master și un server sclave. Cele două servere îmbunătățesc operațiile de citire și scriere și asigură redundanța datelor. \end{itemize}

Semnificația practică a acestei lucrări constă în faptul că rezultatele obținute în urma aplicării soluțiilor scalabile oferă o perspectivă asupra traficului de date și de cereri la care poate rezista o aplicație web și cum ar trebui îmbunătățită și optimizată pentru a gestiona eficient cererile utilizatorilor.

\section{Structura lucrării}
În capitolele ce urmează, se vor prezenta următoarele idei:  aspectele teoretice ale proiectului, tehnologiile folosite în implementare și structura aplicației. Acestea sunt urmate de descrierea arhitecturii, în care sunt evidențiate diferitele module ce constituie aplicația și interacțiunea dintre ele la nivel scalabil. 
Capitolele vor cuprinde și exemplificări cu fragmente de cod care un rol crucial în aplicarea fenomenului de scalabilitate asupra aplicației web.
De asemenea, lucrarea abordează explicații amănunțite despre modul de funcționare și utilizare a interfeței grafice a aplicației cu ajutorul capturilor de ecran.

\chapter{Aspectele teoretice ale proiectului}
În mod ideal, scalabilitatea presupune gestionarea optimă a resurselor hardware, dar și software, astfel încât aplicația când este expusă unui număr mare de cereri le gestionează eficient, fără a compromite viteza de răspuns sau resursele disponibile. Scalabilitatea software \cite{ss} se referă la implementarea unor soluțiilor de tip software care asigură funcționarea eficientă a aplicației chiar și la un număr mare de utilizatori. Aceasta folosită concomitent cu scalabilitatea hardware, care are rol în extinderea resurselor fizice ale unui sistem, aduce la rezultate impresionante în gestionarea traficului și a numărului de cereri mare pentru a oferi utlizatorilor o experiență fluidă.

În practică, oricât de mult am încerca să avem o viteză de răspuns cât mai optimă, aceasta implică un cost foarte mare al resurselor software, hardware dar și a energiei electrice folosite. Un volum limitat de resurse hardware și software implică, de asemenea, o eficiență de gestionare a cererilor dintr-o aplicației mai scăzută. Astfel, este necesară o arhitectură software și hardware bine definită pentru a reuși în atingerea unor rezultate promițătoare cu un număr minim de resurse folosite. Această lucrare urmărește să evidențieze importanța soluțiilor scalabile software implementate în cadrul unei aplicații de tip rezervări web care are rol în optimizarea și fluidizarea volumului de traficul de date rezultat în urma cererilor transmise de utilizatori către server.

\textbf{O aplicație web} \cite{web} este o aplicație de tip software accesată prin intermediul unui browser web care poate rula pe internet sau pe un server. Aceasta poate fi accesată de pe orice tip de dispozitiv conectat la internet, cum ar fi calculatoarele, telefoanele, consolele și este implementată prin intermediul unor limbaje de programare web ca HTML, CSS și JavaScript. Ceea ce face diferită o aplicație web față de cele mobile sau desktop este faptul că poate să ruleze pe unul sau mai multe servere și este accesată cu ajutorul unui browser web. Principalele caracteristici și funcționalități ale unei aplicații web sunt:
\begin{itemize}
\item Scalabilitatea: datorită faptului că aplicațiile web rulează pe un server există posibilitatea de a implementa soluții scalabile care să îmbunătățească și optimizeze performanța pentru a răspunde la cerințele utilizatorilor.
\item Accesibilitatea: oricine are un dispozitiv cu conexiune la internet poate să acceseze aplicațiile web, indiferent de locație sau specificațiile hardware și software pe care le are dispozitivul folosit.
\item Integrarea cu alte servicii și aplicații: aplicații web fiind versatile, programatorii au posibilitatea de a rafina experiența utilizatorilor cu ajutorul serviciilor și aplicațiilor web suplimentare concepute pentru a personaliza experiența fiecărei persoane.
\item Actualizarea eficientă: odată ce s-a confirmat de către echipa de testare faptul că noua versiune a aplicației web a trecut toate testele și îndeplinește toate cerințele, aceasta poate fi lansată către public fără a fi necesară intervenția din partea lui, actualizarea fiind automată pe server.
\end{itemize}

 \section {Protocolul HTTP}
În contextul vremurilor actuale, aplicațiile web au devenit unul dintre cele mai rapide, populare, dar și accesibile metode de a accesa informații și servicii de către utilizatori doar la un click distanță. La baza acestora se află o arhitectură de tip client-server, unde interacțiunele sunt realizate prin intermediul protocoalelor de comunicație.

Protocolul HTTP (Hypertext Transfer Protocol) \cite{http} este protocolul de bază al oricărui schimb de date pe internet care se ocupă cu preluarea resurselor, cum ar fi documentele HTML (HyperText Markup Language), fișiere audio, video ș.a.m.d. Fiind un protocol de tip client-server, cererile sunt inițializate de către client (browser-ul web). Comunicarea dintre clienți și servere are loc prin schimbul de mesaje individuale, unde mesajele trimise de client se numesc cereri, iar mesajele trimise de către server se numesc răspunsuri.

 \subsection{Componentele sistemelor bazate pe HTTP}
HTTP este un protocol de nivel de aplicație ce urmează un model de tip client-server în care un client începe o conexiune pentru a face o cerere, care așteaptă până primește un răspuns de la server. De asemenea, HTTP este un protocol fără reținerea stării, unde serverul nu reține date despre sesiunea activă.

\subsubsection{Clientul: user-agent-ul}
User-agent-ul \cite{usag} este un header de tip string folosit într-o cerere HTTP pentru a trimite informații despre aplicație, sistemul de operare, furnizorul și/sau versiunea folosită către server. 

 \begin{figure}[H]
\centering
\includegraphics*[width=0.8\columnwidth, center]{figuri/ussyntax.png}
\caption{Sintaxa user-agent-ului }
\label{fig:ussyntax}
\end{figure}


Ca să afișeze o pagină web, în prima fază browser-ul trimite o cerere inițială pentru a prelua documentul HTML care reprezintă pagina. În cea de-a doua fază, acesta analizează fișierul primit și trimite cereri suplimentare pentru a obține informații despre stilizarea paginii (CSS), eventuale scripturi ce trebuiesc executate și sub-resursele aflate în pagină (care pot fi imagini sau videoclipuri). În fază finală a procesului, browser-ul combină toate resursele pentru a prezenta documentul completat.
 
\subsection {Server-ul web}
Server-ul web \cite{svweb} se ocupă cu livrarea documentelor cerute de către client. Din punct de vedere al structurii hardware, poate fi un calculator sau o colecție de servere care împart cererile ce au rol în conectarea la internet și la schimbul de date stocate (de exemplu documente HTML, imaginii, fișiere JavaScript) cu alte dispozitive.
Pe de altă parte, din punct de vedere al structurii software, un server web are rolul de a înțelege adresele URL și protocolul HTTP.

URL (Uniform Resource Locators) \cite {url} reprezintă adresa unei resurse unice de pe internet și unul dintre mecanismele folosite de browsere pentru a prelua resurse cum ar fi pagini HTML, documente, imagini ș.a.m.d. Orice URL introdus în bara de adrese duce la o solicitare directă către un server web pentru a accesa resursa specifică, iar în funcție de răspunsul server-ului, resursa solicitată va fi afișată. Un URL este compus din diferiți parametrii, unii fiind obligatorii, altele opționale. 
 \begin{figure}[H]
\centering
\includegraphics*[width=1.7\columnwidth, center]{figuri/url.png}
\caption{Componentele URL-ului}
\label{fig:wbsv}
\end{figure}

\subsubsection{Tipul de protocol URL}
Primul fragment din URL reprezintă tipul de protocol pe care browser-ul trebuie să îl folosească pentru a trimite o cerere de acces la resursă. Protocoalele cele mai întâlnite sunt HTTPS (Hypertext Transfer Protocol Secured) sau HTTP (varianta nesecurizată), dar există cazuri în care este folosit și protocolul mailto: care deschide un client de tip mail.

\subsubsection{Numele resursei și portul}
Numele resursei sau numele domeniului, care este separat de tipul de protocol prin tiparul de caractere \textbf{://}, reprezintă adresa paginii web redactată sub forma unui string. Orice pagină web are atribuit un IP public unic din intervalul 1.0.0.0 - 9.255.255.255, numere care sunt greu de reținut sau scris pentru orice persoană. Din acest motiv, numele domeniului ajută în fluidizarea procesului de accesare a unei pagini web prin „traducerea în limbaj natural” a adreselor IP într-un format ușor de înteles și accesibil pentru oricine.
Folosirea portului este opțională 



Un server HTTP poate fi direct accesat prin numele domeniilor pe care site-urile web le dețin și livrează conținutul acestora către utilizatori.
La nivel fundamental, modul de funcționare al unui server web presupune ca de fiecare dată când un browser are nevoie de un fișier găzduit de către un server web, acesta va trimite o cerere HTTP pentru a primi fișierul respectiv. Odată ce server-ul HTTP a recepționat cererea, verifică dacă URL solicitat există și în cazul confirmării, acesta trimite la utilizator conținutul cerut. În cazul în care server-ul nu reușește să verifice cererea, mai întâi acesta verifică dacă este necesară generarea unui fișier dinamic. În cele din urmă, dacă niciuna dintre opțiuni nu este disponibilă, server-ul web va returna codul de eroare 404 Not Found.

 \begin{figure}[H]
\centering
\includegraphics*[width=0.6\columnwidth, center]{figuri/wbsv.png}
\caption{Diagrama modului de funcționare a server-ului web}
\label{fig:wbsv}
\end{figure}

\subsection{Cum funcționează o conexiune HTTP}
O conexiune HTTP este controlată la nivelul stratului de transport, ceea ce înseamnă că, de fapt, nu este rolul protocolului HTTP de a-și gestiona conexiunea, ci de a se folosi de un protocol de transport. Acesta are nevoie de un protocol de transport fiabil, prin care mesajele au o rată de pierdere foarte mică. Din acest motiv, se utilizează protocolul TCP datorită fiabilității sale care garantează să trimită mesajele cu un număr minim erori, deși acesta are o viteză de răspuns mai lentă în comparație cu protocolul UDP. 
 \begin{figure}[H]
\centering
\includegraphics*[width=0.8\columnwidth, center]{figuri/tcp.png}
\caption{Diagrama conexiunii HTTP}
\label{fig:tcp}
\end{figure}
Primul pas în a începe o conexiune HTTP o reprezintă schimbul de cerere/răspuns dintre client și server. Acest schimb începe odată ce s-a stabilit o conexiune TCP, un proces care trimite unul sau mai multe cereri din partea clientului pentru a primi un răspuns de la server. Clientul are posibilitatea de a deschide o nouă conexiune, de a folosi o conexiune existentă sau de a deschide mai multe conexiuni TCP către servere.
Al doilea pas presupune trimiterea mesajul HTTP către server și răspunsul server-ului la cererea clientului cu informația cerută. Dacă răspunsul serverului are codul 200 OK, înseamnă faptul că s-a stabilit cu succes conexiunea, iar serverul poate răspunde clientului cu informația cerută.  
 \begin{figure}[H]
\centering
\includegraphics*[width=0.7\columnwidth, center]{figuri/cltcp.png}
\caption{Cerere HTTP}
\label{fig:tcp}
\end{figure}

 \begin{figure}[H]
\centering
\includegraphics*[width=0.8\columnwidth, center]{figuri/svtcp.png}
\caption{Răspuns HTTP}
\label{fig:tcp}
\end{figure}
 
În cele din urmă, conexiunea este închisă sau reutilizată.

\subsection{Mesaje HTTP}


 Alegerea de a utiliza acest protocol în cadrul dezvoltării aplicației este justificată prin simplitatea cu care se implementează acesta, dar și interoperabilitatea sa extinsă cu browserele, serverele web și backend-urile.

 

\chapter{Structura aplicației}
Aplicația a fost structurată în șapte module, fiecare modul având un rol bine definit în arhitectura generală. Acestea sunt următoarele:
\begin{enumerate}
\item Modulul de aplicație web în Python Flask este prezent pe stația de lucru și include rute importante precum /, /login, /register, /logout, /profile, /admin, /reserve pentru gestionarea logicii unui sistem de management al rezervărilor, autentificarea utilizatorilor și interacțiunea cu baza de date. Funcționalitățile rutelor sunt:
\begin{itemize}
\item /login - autentificarea utilizatorilor și validarea credențialelor.
\item /register - înregistrarea utilizatorilor noi și validarea datelor.
\item /logout - finalizarea sesiunii active.
\item /reserve - rezervarea unui hotel (operațiune de tip CRUD).
\item /profile - afișarea și actualizarea profilelor utilizatorilor.
\item /admin - rută pentru utilizatorii cu rol de administrator.
\item / - pagina de acasă cu o interfață generală sau redirecționare în funcție de sesiune.
\end{itemize}
\item Modulul de load balancing (Nginx): Nginx (instalat pe stația de lucru) se ocupă de distribuirea cererilor primite de la utilizatori către cele două servere: aplicația locală Flask și serverul IIS de pe VM2, ceea ce îmbunătățește scalabilitatea soluției implementate. Load balancer-ul are mai multe roluri cum ar fi:
\begin{itemize}
\item Reverse proxy care asigură redirecționarea traficului.
\item Verificarea și detectarea automată a sănătății server-ului și scoaterea temporară din funcționare până când sunt rezolvate erorile.
\item Balansarea cererilor trimise către servere, astfel încât cereri sunt preluate secvențial.
\item Măsură de rezervă în cazul în care unul dintre servere devine inaccesibil, traficul va fi redirecționat către următorul server disponibil.
\end{itemize}
\item Modulul de testare are rol în testarea izolată a aplicației web fără ca sistemul principal să fie afectat. Pentru acest mediu de testare izolat s-a folosit VM1 care are responsabilitatea de a testa codul Flask, conectivitatea cu serviciul MySQL, validarea configurației Nginx, rularea unui debug server Flask și de utilizarea a aplicației Postman sau chiar a browser-ului local pentru a rula testele.
\item Modulul de server IIS reprezintă o a doua instanță web care poate fi folosit pentru a replica parțial funcționalități ale aplicației Flask, distribuirea paginilor, resurselor statice (HTML, JavaScript, CSS), resurselor REST și pentru a răspunde la cererile distribuite de către load balancer.
\item Modulul de bază de date MySQL Master, ce se află pe stația de lucru, are la bază nodul principal al bazei de date unde toate operațiile de scriere (INSERT, UPDATE, DELETE) se fac pe acest server. Principalele obiective ale modulului sunt de a gestiona tabelele aplicației web, replicarea automată către slave (VM2), creare utilizatorilor și a privilegiilor pentru replicare și acceptarea conexiunilor de la aplicația Flask.
\item Modulul de bază de date MySQL Sclave are rolul de a prelua și stoca copii ale datelor scrise pe master în timp real pentru a efectua operații de citire (SELECT) asupra bazei de date. Această tehnică asigură îmbunătățirea performanței aplicației atunci când este folosită frecvent comanda SELECT și asigură redundanța datelor. Un alt rol crucial pe care îl reprezintă sclave este acela că poate fi promovat la master în cazul în care nodul principal eșuează, astfel datele nu vor putea fi pierdute sau compromise.
\item Modulul de rețea și virtualizare asigură comunicarea între stația de lucru, VM1 și VM2 prin intermediul rețelei bridged sau NAT în VMware Workstation. La nivelul fiecărei mașini sunt setate IP-uri statice pentru a asigura o conexiune stabilă.
\end{enumerate}



\chapter{Noțiuni teoretice privind aplicațiile web și scalabilitatea}

   Pentru facilitarea formatării tezei de licență/dizertație în
   \LaTeX{}, se furnizează fișierul clasă ``upb\_thesis.cls'' și un model de
   folosire ``exemplu\_teza.tex''.  Utilizarea comenzilor standard
   \verb=\documentclass=, \verb=usepackage=, \verb=setcounter=,
   etc. nu este diferită de utilizarea lor în orice alt context
   \LaTeX{}. Comenzile standard sau cele specifice {UPB}
   apar într-o secvență care trebuie respectată pentru a obține documentul în formatul
   acceptat de universitate.  Pentru redactarea tezei, este
   obligatorie folosirea următoarelor comenzi:

\verb=\singlespacing= specifică spațierea pentru întregul document.

\verb=\author{Ion Creangă}= numele studentului va fi inserat pe pagina de titlu și în declarația de onestitate.  

\verb=\title{Basmele românilor}= titlul tezei 

\verb=\facultatea{Facultatea de Electronică, Telecomunicații și Tehnologia= \\
\verb=Informației}=

\verb=\tiplucrare{licență}= parametrul trebuie să fie ``licență'' sau ``dizertație''

\verb=\domeniu{Electronică, Telecomunicații și Tehnologia Informației}=

\verb=\catedra{}=Conform anexei 4 din ghidula absolventului. 

\verb=\program{Povești}= numele specializării

\verb=\titlulobtinut{Inginer}= parametrul trebuie să fie ``Inginer'' sau ``Master''

\verb=\director{Mihail Kogălniceanu}= numele îndrumătorului de
proiect. Dacă sunt mai mulți, se vor separate cu doi backslash și un spațiu: 

\verb=\director{Mihail Kogălniceanu\\ Titu Maiorescu}=

\verb=\submissionmonth{Iunie}=

\verb=\submissionyear{2011}= data care apare pe pagina de titlu

Listele opționale de figuri, tabele, și abrevieri trebuie incluse în secvența \\
\verb=\beforepreface= - \verb=\afterpreface=: 
\begin{verbatim}
\beforepreface
\listoffigures
\listoftables
\abbreviations{
}
\afterpreface
\end{verbatim}
\section{Alte recomandări}
\begin{itemize}
\item Se recomandă folosirea pachetului aspell (pachetul \verb=aspell-ro= în
Debian/Ubuntu) pentru verificarea ortografiei, cu comanda:
{\small \begin{verbatim}aspell --lang=ro -t check  ./exemplu_teza.tex\end{verbatim}}

\item Se recomandă verificarea fonturilor în PDF-ul produs - este preferabil
să nu se folosească decât ``Type 1'' sau ``Truetype'', nu ``Type 3''. 
\end{itemize}

\chapter{Basme}
\section{Punguța cu doi bani}


Era odată o babă și un moșneag. Baba avea o găină, și
moșneagul un cucoș; găina babei se oua de câte două ori pe fiecare zi
și baba mânca o mulțime de ouă; iar moșneagului nu-i da nici
unul. Moșneagul într-o zi perdu răbdarea și zise:

— Măi babă, mănânci ca în târgul lui Cremene. Ia dă-mi și mie niște
ouă, ca să-mi prind pofta măcar. — Da' cum nu! zise baba, care era
foarte zgârcită. Dacă ai poftă de ouă, bate și tu cucoșul tău, să facă
ouă, și-i mânca; că eu așa am bătut găina, și iacătă-o cum se ouă.

Moșneagul, pofticios și hapsin, se ia după gura babei și, de ciudă,
prinde iute și degrabă cucoșul și-i dă o bataie bună, zicând:

— Na! ori te ouă, ori du-te de la casa mea; ca să nu mai strici
mâncarea degeaba.

Cucoșul, cum scăpă din mânile moșneagului, fugi de-acasă și umbla pe
drumuri, bezmetec. Și cum mergea el pe-un drum, numai iată găsește o
punguță cu doi bani (Figura \ref{fig:punguta}). Și cum o găsește, o și ia în clonț și se întoarnă
cu dânsa înapoi către casa moșneagului. Pe drum se întâlnește c-o
trăsură c-un boier și cu niște cucoane. Boierul se uită cu băgare de
seamă la cucoș, vede în clonțu-i o punguță și zice vezeteului:

— Măi! ia dă-te jos și vezi ce are cucoșul cela în plisc.

\begin{figure}[tb]
\centering
\includegraphics*[width=0.4\columnwidth]{figuri/punguta.jpg}
\caption{Cucoșu găsește o pungă cu doi galbeni.}
\label{fig:punguta}
\end{figure}



Vezeteul se dă iute jos din capra trăsurei, și c-un feliu de meșteșug,
prinde cucoșul și luându-i punguța din clonț o dă boieriului. Boieriul
o ia, fără păsare o pune în buzunar și pornește cu trăsura
înainte. Cucoșul, supărat de asta, nu se lasă, ci se ia după trăsură,
spuind neîncetat:

Cucurigu ! boieri mari, Dați punguța cu doi bani !


[...]

(Publicată pentru prima oară în \cite{punguta1876}.)

\section{Tinerețe fără bătrânețe și viață fără de moarte}

A fost odată ca niciodată; că de n-ar fi, nu s-ar mai povesti; de când
făcea plopșorul pere și răchita micșunele; de când se băteau urșii în
coade; de când se luau de gât lupii cu mieii de se sărutau,
înfrățindu-se; de când se potcovea puricele la un picior cu nouăzeci
și nouă de oca de fier și s-arunca în slava cerului de ne aducea
povești;


    De când se scria musca pe părete,
    Mai mincinos cine nu crede.


A fost odată un împărat mare și o împărăteasă, amândoi tineri și
frumoși, și, voind să aibă copii, a făcut de mai multe ori tot ce
trebuia să facă pentru aceasta; a îmblat pe la vraci\cite{TanenbaumCN02} și filosofi, ca
să caute la stele și să le ghicească daca or să facă copii; dar în
zadar. În sfârșit, auzind împăratul că este la un sat, aproape, un
unchiaș dibaci, a trimis să-l cheme; dar el răspunse trimișilor că:
cine are trebuință, să vie la dânsul. S-au sculat deci împăratul și
împărăteasa și, luând cu dânșii vro câțiva boieri mari, ostași și
slujitori, s-au dus la unchiaș acasă. Unchiașul, cum i-a văzut de
departe, a ieșit să-i întâmpine și totodată le-a zis:

- Bine ați venit sănătoși; dar ce îmbli, împărate, să afli? Dorința ce
ai o să-ți aducă întristare.

- Eu nu am venit să te întreb asta, zise împăratul, ci, daca ai ceva
leacuri care să ne facă să avem copii, să-mi dai.

- Am, răspunse unchiașul; dar numai un copil o să faceți. El o să fie
Făt-Frumos și drăgăstos, și parte n-o să aveți de el. Luând împăratul
și împărăteasa leacurile, s-au întors veseli la palat și peste câteva
zile împărăteasa s-a simțit însărcinată. Toată împărăția și toată
curtea și toți slujitorii s-au veselit de această întâmplare.

Mai-nainte de a veni ceasul nașterii, copilul se puse pe un plâns, de
n-a putut nici un vraci să-l împace. Atunci împăratul a început să-i
făgăduiască toate bunurile din lume, dar nici așa n-a fost cu putință
să-l facă să tacă.

- Taci, dragul tatei, zice împăratul, că ți-oi da împărăția cutare sau
cutare; taci, fiule, că ți-oi da soție pe cutare sau cutare fată de
împărat, și alte multe d-alde astea; în sfârșit, dacă văzu și văzu că
nu tace, îi mai zise: taci, fătul meu, că ți-oi da Tinerețe fără
bătrânețe și viață fără de moarte.

Atunci, copilul tăcu și se născu; iar slujitorii deteră în timpine și
în surle și în toată împărăția se ținu veselie mare o săptămână
întreagă.

De ce creștea copilul, d-aceea se făcea mai isteț și mai îndrăzneț. Îl
deteră pe la școli și filosofi, și toate învățăturile pe care alți
copii le învăța într-un an, el le învăța într-o lună, astfel încât
împăratul murea și învia de bucurie. Toată împărăția se fălea că o să
aibă un împărat înțelept și procopsit ca Solomon împărat. De la o
vreme încoace însă, nu știu ce avea, că era tot galeș, trist și dus pe
gânduri. Iar când fuse într-o zi, tocmai când copilul împlinea
cincisprezece ani și împăratul se afla la masă cu toți boierii și
slujbașii împărăției și se chefuiau, se sculă Făt-Frumos și zise:

- Tată, a venit vremea să-mi dai ceea ce mi-ai făgăduit la naștere.

Auzind aceasta, împăratul s-a întristat foarte și i-a zis:

– Dar bine, fiule, de unde pot eu să-ți dau un astfel de lucru nemaiauzit? Și dacă ți-am făgăduit atunci, a fost numai ca să te împac.

– Daca tu, tată, nu poți să-mi dai, apoi sunt nevoit să cutreier toată lumea până ce voi găsi făgăduința pentru care m-am născut.

Atunci toți boierii și împăratul deteră în genuchi, cu rugăciune să nu părăsească împărăția; fiindcă, ziceau boierii:

– Tatăl tău de aci înainte e bătrân, și o să te ridicăm pe tine în scaun, și avem să-ți aducem cea mai frumoasă împărăteasă de sub soare de soție.

Dar n-a fost putință să-l întoarcă din hotărârea sa, rămânând statornic ca o piatră în vorbele lui; iar tată-său, dacă văzu și văzu, îi dete voie și puse la cale să-i gătească de drum merinde și tot ce-i trebuia.


[...]

\section[Cinci Pâini]{Cinci pâini \footnote{Anecdotă publicată prima oară în Convorbiri literare, nr. 12, 1 martie 1883}}


Doi oameni, cunoscuți unul cu altul, călătoreau odată, vara, pe un drum. Unul avea în traista sa trei pâni, și celalalt două pâni. De la o vreme, fiindu-le foame, poposesc la umbra unei răchiți pletoase, lângă o fântână cu ciutură, scoate fiecare pânile ce avea și se pun să mănânce împreună, ca să aibă mai mare poftă de mâncare.


Tocmai când scoaseră pânile din traiste, iaca un al treile drumeț, necunoscut, îi ajunge din urmă și se oprește lângă dânșii, dându-le ziua bună. Apoi se roagă să-i deie și lui ceva de mâncare, căci e tare flămând și n-are nimica merinde la dânsul, nici de unde cumpăra.

— Poftim, om bun, de-i ospăta împreună cu noi, ziseră cei doi drumeți călătorului străin; căci mila Domnului! unde mănâncă doi mai poate mânca și al treilea.

Călătorul străin, flămând cum era, nemaiașteptând multă poftire, se așază jos lângă cei doi, și încep a mânca cu toții pâne goală și a be apă rece din fântână, căci altă udătură nu aveau. Și mănâncă ei la un loc tustrei, și mănâncă, până ce gătesc de mâncat toate cele cinci pâni, de parcă n-au mai fost.

După ce-au mântuit de mâncat, călătorul străin scoate cinci lei din pungă și-i dă, din întâmplare, celui ce avusese trei pâni, zicând:

— Primiți, vă rog, oameni buni, această mică mulțămită de la mine, pentru că mi-ați dat demâncare la nevoie; veți cinsti mai încolo câte un pahar de vin, sau veți face cu banii ce veți pofti. Nu sunt vrednic să vă mulțămesc de binele ce mi-ați făcut, căci nu vedeam lumea înaintea ochilor de flămând ce eram.

Cei doi nu prea voiau să primească, dar, după multă stăruință din partea celui al treilea, au primit. De la o vreme, călătorul străin și-a luat ziua bună de la cei doi și apoi și-a căutat de drum. Ceilalți mai rămân oleacă sub răchită, la umbră, să odihnească bucatele. Și, din vorbă în vorbă, cel ce avuse trei pâni dă doi lei celui cu două pâni, zicând:

— Ține, frate, partea dumitale, și fă ce vrei cu dânsa. Ai avut două pâni întregi, doi lei ți se cuvin. Și mie îmi opresc trei lei, fiindc-am avut trei pâni întregi, și tot ca ale tale de mari, după cum știi.

— Cum așa?! zise celălalt cu dispreț! pentru ce numai doi lei, și nu doi și jumătate, partea dreaptă ce ni se cuvine fiecăruia? Omul putea să nu ne deie nimic, și atunci cum rămânea?

— Cum să rămâie? zise cel cu trei pâni; atunci aș fi avut eu pomană pentru partea ce mi se cuvine de la trei pâni, iar tu, de la două, și pace bună. Acum, însă, noi am mâncat degeaba, și banii pentru pâne îi avem în pungă cu prisos: eu trei lei și tu doi lei, fiecare după numărul pânilor ce am avut. Mai dreaptă împărțeală decât aceasta nu cred că se mai poate nici la Dumnezeu sfântul...

— Ba nu, prietene, zice cel cu două pâni. Eu nu mă țin că mi-ai făcut parte dreaptă. Haide să ne judecăm, și cum a zice judecata, așa să rămâie.

— Haide și la judecată, zise celălalt, dacă nu te mulțămești. Cred că și judecata are să-mi găsească dreptate, deși nu m-am târât prin judecăți de când sunt.

Și așa, pornesc ei la drum, cu hotărârea să se judece. Și cum ajung într-un loc unde era judecătorie, se înfățoșează înaintea judecătorului și încep a spune împrejurarea din capăt, pe rând fiecare; cum a venit întâmplarea de au călătorit împreună, de au stat la masă împreună, câte pâni a avut fiecare, cum a mâncat drumețul cel străin la masa lor, deopotrivă cu dânșii, cum le-a dat cinci lei drept mulțămită și cum cel cu trei pâni a găsit cu cale să-i împartă.

Judecătorul, după ce-i ascultă pe amândoi cu luare aminte, zise celui cu două pâni: — Și nu ești mulțămit cu împărțeala ce s-a făcut, omule?

— Nu, domnule judecător, zise nemulțămitul; noi n-am avut de gând să luăm plată de la drumețul străin pentru mâncarea ce i-am dat; dar, dac-a venit întâmplarea de-așa, apoi trebuie să împărțim drept în două ceea ce ne-a dăruit oaspetele nostru. Așa cred eu că ar fi cu cale, când e vorba de dreptate.

— Dacă e vorba de dreptate, zise judecătorul, apoi fă bine de înapoiește un leu istuialalt, care spui c-a avut trei pâni.

— De asta chiar mă cuprinde mirare, domnule judecător, zise nemulțămitul cu îndrăzneală. Eu am venit înaintea judecăței să capăt dreptate, și văd că dumneata, care știi legile, mai rău mă acufunzi. De-a fi să fie tot așa și judecata dinaintea lui Dumnezeu, apoi vai de lume!

— Așa ți se pare dumitale, zise judecătorul liniștit, dar ia să vezi că nu-i așa. Ai avut dumneata două pâni?

— Da, domnule judecător, două am avut.

— Tovarășul dumitale, avut-a trei pâni?

— Da, domnule judecător, trei a avut.

— Udătură ceva avut-ați vreunul?

— Nimic, domnule judecător, numai pâne goală și apă răce din fântână, fie de sufletul cui a făcut-o acolo, în calea trecătorilor.

— Dinioarea, parcă singur mi-ai spus, zise judecătorul, că ați mâncat toți tot ca unul de mult; așa este?

— Așa este domnule judecător.

— Acum, ia să statornicim rânduiala următoare, ca să se poată ști hotărât care câtă pâne a mâncat. Să zicem că s-a tăiat fiecare pâne în câte trei bucăți deopotrivă de mari; câte bucăți ai fi avut dumneata, care spui că avuși două pâni?

— Șese bucăți aș fi avut, domnule judecător.

— Dar tovarășul dumitale, care spui că avu trei pâni?

— Nouă bucăți ar fi avut, domnule judecător.

— Acum, câte fac la un loc șese bucăți și cu nouă bucăți?

— Cincisprezece bucăți, domnule judecător.

— Câți oameni ați mâncat aceste cincisprezece bucăți de pâne?

— Trei oameni, domnule judecător.

— Bun! Câte câte bucăți vin de fiecare om?

— Câte cinci bucăți, domnule judecător.

— Acum, ții minte câte bucăți ai fi avut dumneta?

— Șese bucăți, domnule judecător.

— Dar de mâncat, câte ai mâncat dumneta?

— Cinci bucăți, domnule judecător.

— Și câte ți-au mai rămas de întrecut?

— Numai o bucată, domnule judecător.

— Acum să stăm aici, în ceea ce te privește pe dumneta, și să luăm pe istalalt la rând. Ții minte câte bucăți de pâne ar fi avut tovarășul d-tale?

— Nouă bucăți, domnule judecător.

— Și câte a mâncat el de toate?

— Cinci bucăți, ca și mine, domnule judecător.

— Dar de întrecut, câte i-au mai rămas?

— Patru bucăți, domnule judecător.

— Bun! Ia, acuș avem să ne înțelegem cât se poate de bine! Vra să zică, dumneta ai avut numai o bucată de întrecut, iar tovarășul dumitale, patru bucăți. Acum, o bucată de pâne rămasă de la dumneta și cu patru bucăți de la istalalt fac la un loc cinci bucăți?

— Taman cinci, domnule judecător.

— Este adevărat că aceste bucăți de pâne le-a mâncat oaspetele dumneavoastră, care spui că v-a dat cinci lei drept mulțămită?

— Adevărat este, domnule judecător.

— Așadar, dumitale ți se cuvine numai un leu, fiindcă numai o bucată de pâne ai avut de întrecut, și aceasta ca și cum ai fi avut-o de vânzare, deoarece ați 
primit bani de la oaspetele dumneavoastră. Iar tovarășul dumitale i se cuvin patru lei, fiindcă patru bucăți de pâne a avut de întrecut. Acum, dară, fă bine de înapoiește un leu tovarășului dumitale. Și dacă te crezi nedreptățit, du-te și la Dumnezeu, și las' dacă ți-a face și el judecată mai dreaptă decât aceasta!

Cel cu două pâni, văzând că nu mai are încotro șovăi, înapoiește un leu tovarășului său, cam cu părere de rău, și pleacă rușinat. Cel cu trei pâni însă, uimit de așa judecată, mulțămește judecătorului și apoi iese, zicând cu mirare:

— Dac-ar fi pretutindene tot asemenea judecători, ce nu iubesc a li cânta cucul din față, cei ce n-au dreptate n-ar mai năzui în veci și-n pururea la judecată.

Corciogarii, porecliți și apărători, nemaiavând chip de traiu numai din minciuni, sau s-ar apuca de muncă, sau ar trebui, în toată viața lor, să tragă pe dracul de coadă... Iar societatea bună ar rămâne nebântuită.



\bibliographystyle{unsrt}
\bibliography{referinte}


\appendix
\chapter {Rezolvare ''Cinci Pâini''}
A aduce 3 pâini, iar B aduce 2 pâini. C plătește 5 lei.  Se împarte
fiecare pâine în 3 părți egale: $3 \times 3 + 2 \times 3 = 15$
bucăți. Fiecare mesean capătă câte 5 bucăți. Reiese că A i-a dat lui C
$3 \times 3 - 5 = 4$ pâini, iar B i-a dat $2 \times 3 - 5 = 1$
pâine. Împărțeală cinstită este deci 4 lei pentru A, și 1 leu pentru B.

\chapter{Alte povești și basme}

 Vezi tabela \ref{altebasme} cu alte povești și autorii/culegătorii lor.
\lstinputlisting[language=Matlab]{Cod/Cod.m}


\begin{table}
\begin{tabular}[t]{c|p{3cm}|p{3cm}|p{3cm}|c}
\_ & Dănilă prepeleac & Prâslea cel voinic și merele de aur & Făt-Frumos din lacrimă & Greuceanu \\
\hline
 & Ion Creangă & Petre Ispirescu & Mihai Eminescu & Petre Ispirescu \\
\end{tabular}
\label{altebasme}
\caption{Alte basme.}
\end{table}





\end{document}
